\section{Illumination and Speech}\label{sec:illumination}

God governs the world through the angels, and he first governs the angels
through one another. Having shown that God moves every creature and works
in every agent (\ST{I qq.~103--105}), Aquinas turns to the way ``one
creature moves another'' (\latin{quomodo una creatura moveat aliam}) and
divides the consideration threefold: how the angels move, \latin{qui sunt
creaturae pure spirituales}; how bodies move; and how men move, who are
composed of both natures. Under the angels he asks first how angel acts on
angel, then on the corporeal creature (q.~110), then on men (q.~111); and
under angel acting on angel he treats their enlightenment, their speech,
and their ordering to one another, good and bad alike (\latin{de
illuminatione, et locutione Angelorum, et ordinatione eorum ad invicem, tam
bonorum, quam malorum}) (\ST{I q.~106 pr.}). Illumination and speech (qq.~106--107) belong here; the ordering
of the good and the bad angels (qq.~108--109) follows in
\S\ref{sec:hierarchies}; the government of bodies and of men follows in
\S\ref{sec:government}.

The governing principle is Aquinas's own answer to the question whether
God governs all things immediately. In the design of government he does;
in its execution ``He governs some things by means of others.'' For ``it is
a greater perfection for a thing to be good in itself and also the cause of
goodness in others, than only to be good in itself. Therefore God so
governs things that He makes some of them to be causes of others in
government; as a master, who not only imparts knowledge to his pupils, but
gives also the faculty of teaching others'' (\ST{I q.~103 a.~6}). The
\emph{sed contra} of that article is Augustine's picture of the whole
creation ruled in order: ``the sinful and unfaithful spirit is ruled by the
good and just spirit of life; and this spirit by God Himself'' (as quoted
in \ST{I q.~103 a.~6 s.c.}; Augustine, \emph{De Trinitate} III.4.9).
Dionysius calls the highest dignity of any member of a hierarchy the dignity
``of becoming a fellow-worker with God'' (\emph{CH} 3.2), and Aquinas, in
the disputed questions on truth, makes that sentence the root of angelic
illumination. From the divine goodness it proceeds that God shares his
perfection with creatures in proportion to them, not only so that they are
good and perfect in themselves, but so that they bestow perfection on
others, \latin{Deo quodammodo cooperando}. \latin{Et hic est nobilissimus
modus divinae imitationis}: this is the noblest way of imitating God, and
from it proceeds the order among the angels by which some enlighten others
(\emph{De veritate} q.~9 a.~2). The \emph{Divine
Names} gives the same doctrine its most luminous form. The intelligible
essences ``are illuminated as to the reasons of things, in a manner
peculiar to themselves; and they again convey to their kindred spirits
things appropriate to them''; from the Good come to them ``the powers
elevating the inferior to the superior, the providences of the more exalted
for those below them,'' so ``that angels are, as it were, heralds of the
Divine silence, and project, as it were, luminous lights revealing Him Who
is in secret'' (\emph{DN} 4.1--2).

Scripture shows the angels speaking to one another and sending one
another. Zechariah hears the angel of the Lord intercede, ``O Lord of
hosts, how long wilt thou not have mercy on Jerusalem,'' and the Lord
answer ``the angel, that spoke in me, good words, comfortable words''
(Zach 1:12--13); then ``the angel that spoke in me went forth, and another
angel went out to meet him. And he said to him: Run, speak to this young
man'' (Zach 2:3--4 [2:7--8]). The Seraphim ``cried one to another'' (Isa 6:3).
Gregory the Great draws the law from Zechariah: \latin{Hoc tamen
certissime scimus, quia ad explendum de supernis ministerium alii spiritus
alios mittunt} --- this we know most certainly, that to fulfil a ministry
from on high some spirits send others; \latin{Dum enim angelus ad angelum
dicit: Curre et loquere, dubium non est, quia alius alium mittit. Minora
vero sunt quae mittuntur, majora quae mittunt} --- the lesser are sent,
the greater send (Gregory, \emph{Hom.\ in Evang.} 34.13). Isidore reads
the same verse to the same effect: were it not that the higher powers
dispose the lower in the very offices of the angels, the angel would never
have learned from an angel what he was to say to a man (\latin{nullo modo
hoc, quod homini diceret angelus, ab angelo cognovisset}; Isidore,
\emph{Etymologiae} VII.5.8).

\subsection{How One Angel Moves Another (\ST{I q.~106})}\label{sec:q106}

The four articles ask whether one angel moves the intellect of another by
enlightening it; whether one moves the will of another; whether a lower
angel can enlighten a higher; and whether the higher enlightens the lower
concerning all that he knows (\ST{I q.~106 pr.}).

\paragraph{One angel enlightens another (a.~1).} The \emph{sed contra} is
Dionysius: ``the angels of the second hierarchy are cleansed, enlightened
and perfected by the angels of the first hierarchy'' (\latin{Angeli
secundae hierarchiae purgantur et illuminantur et perficiuntur per Angelos
primae hierarchiae}; \emph{CH} 8). In Parker's rendering the middle order
``is purified and illuminated and perfected in the manner described, by
the Divine illuminations vouchsafed to it at second hand, through the first
Hierarchical Order'' (\emph{CH} 8.1). Aquinas begins from the nature of
light. ``Intellectual light is nothing else than a manifestation of truth,
according to Ephesians 5:13: `All that is made manifest is light.' Hence to
enlighten means nothing else but to communicate to others the manifestation
of the known truth,'' as the Apostle received the grace ``to enlighten all
men, that they may see what is the dispensation of the mystery which hath
been hidden from eternity in God'' (Eph 3:9). Dionysius says that ``the
subordinate Orders of the Heavenly Beings are taught by the superior, in due
order, the deifying sciences'' (\emph{CH} 7.3).

Since two things concur in understanding, the intellectual power and the
likeness of the thing understood, one angel can make truth known to another
in both. First, ``by strengthening his intellectual power'': as the less hot
grows hotter in the presence of the hotter, ``so the intellectual power of
an inferior angel is strengthened by the superior angel turning to him:
since in spiritual things, for one thing to turn to another, corresponds to
neighborhood in corporeal things'' (\latin{Hoc enim facit in spiritualibus
ordo conversionis, quod facit in corporalibus ordo localis propinquitatis}).
Secondly, on the side of the likeness understood. ``The superior angel
receives the knowledge of truth by a kind of universal conception, to
receive which the inferior angel's intellect is not sufficiently powerful,''
since it is natural to the lower to receive truth more particularly. The
higher therefore ``distinguishes, in a way, the truth which he conceives
universally, so that it can be grasped by the inferior angel,'' as teachers
among us divide into many points what they grasp in the whole,
\latin{providentes capacitati aliorum}
(\ST{I q.~106 a.~1 co.}). Dionysius says it in his exposition of the
teeth under which Scripture pictures the angels: ``The teeth denote the
dividing of the nourishing perfection given to us; for each intellectual
Being divides and multiplies, by a provident faculty, the unified conception
given to it by the more Divine for the proportionate elevation of the
inferior'' (\emph{CH} 15.3).

The replies protect the immediacy of the vision of God. The prophet's
``They shall teach no more every man his neighbour \dots\ for all shall know
me'' (Jer 31:34) holds of the divine essence, which all the blessed angels
see immediately; ``in this respect one does not teach another.'' But of the
types of the divine works (\latin{rationes divinorum operum}), which are
known in God as in their cause, ``each one knows the more types, the more
perfectly he sees God,'' and concerning these the higher enlightens the
lower (ad~1; \emph{DN} 4.1). An angel does not give another the light of
nature, grace, or glory, all of which come from God, but strengthens his
natural light and manifests to him the truth ``concerning the state of
nature, of grace, and of glory'' (ad~2). Augustine taught that the rational
mind is formed by God alone: man is made to the image \latin{ut nulla
natura interposita formetur}, so that no nature intervenes in his forming,
and the spirit made to the image \latin{haeret enim veritati nulla
interposita creatura}, cleaves to the truth with no creature between
(\emph{De diversis quaestionibus LXXXIII} q.~51.2,~4). Aquinas keeps the
teaching whole. The mind is formed immediately by God, both as an image
from its exemplar and as a subject by its ultimate completing form, ``for
the created mind is always considered to be unformed, except it adhere to
the first truth; while the other kinds of enlightenment that proceed from
man or angel, are, as it were, dispositions to this ultimate form'' (ad~3).

The disputed questions on truth set the same doctrine in the analogy of
bodily light, which serves sight in two ways: it makes the potentially
visible actually visible, and it strengthens the eye. Intellectual light
may therefore be called either the vigour of the intellect itself or that by
which something is made known, and one mind enlightens another by handing
it a medium of knowing by which it is strengthened to see what it could not
see before. Among men this happens by speech, as the master enlightens the
disciple, or by sensible signs, as the priest, according to Dionysius, enlightens
the people by ministering and showing them the sacraments, \latin{quae sunt
manuductiones in divina intelligibilia}, guidings by the hand into divine
and intelligible things. The angels come to divine things by neither way:
\latin{Nihil ergo est aliud Angelum ab Angelo illuminari, quam confortari
intellectum inferioris Angeli per aliquid inspectum in superiori, ad aliqua
cognoscenda} --- for an angel to be enlightened by an angel is nothing else
than for the intellect of the lower to be strengthened, by something beheld
in the higher, to know certain things. As in bodies the higher stand to the
lower as act to potency, \latin{ita et superiores spiritus sunt quasi
actus respectu inferiorum}, and every potency is strengthened by its conjunction with its
act; the lower angels are strengthened by their continuity with the
higher, \latin{quae quidem continuatio est per intuitum intellectus},
a continuity made by the gaze of the intellect (\emph{De veritate} q.~9
a.~1). The commentary on the \emph{Sentences} gives the reason the angels
can advance at all: in the vision of God's essence his effects are not all
known of necessity unless his whole power is comprehended, and \latin{quanto
aliquis intellectus limpidius eam videt, tanto plura in ea cognoscere
potest}; the higher therefore see more effects in the divine essence,
\latin{et de illis superiores inferiores illuminare et instruere possunt}
(\emph{In II Sent.} d.~11 q.~2 a.~2).

\paragraph{No angel changes the will of another (a.~2).} The \emph{sed
contra} is an argument from justification: ``To him it belongs to change the
will, to whom it belongs to bestow righteousness: for righteousness is the
rightness of the will. But God alone bestows righteousness.'' The will is
moved in two ways, on the side of its object and on the side of the power.
On the side of the object, the good moves it as the desirable moves desire,
and so does anyone who shows that something is good. Other goods incline
the will somewhat, ``yet nothing sufficiently moves the will save the
universal good, and that is God. And this good He alone shows, that it may
be seen by the blessed,'' as he answered Moses, who asked ``Shew me thy
glory'': ``I will shew thee all good'' (Exod 33:18--19). An angel therefore
moves the will neither as its sufficient object nor as the one who shows it;
``but he inclines the will as something lovable, and as manifesting some
created good ordered to God's goodness. And thus he can incline the will to
the love of the creature or of God, by way of persuasion'' (\latin{per
modum suadentis}). On the side of the power the will can be moved by God
alone, for willing is an inclination of the willer to the thing willed, and
``He alone can change this inclination, Who bestowed on the creature the
power to will,'' since ``He alone is the author of the intellectual nature''
(\ST{I q.~106 a.~2 co.}).

The replies fix the sense of the Dionysian triad. Cleansing, enlightening,
and perfecting are to be understood ``according to the mode of
enlightenment'': God, who enlightens by changing intellect and will, cleanses
both from their defects and perfects both unto their end; but an angel's
enlightenment concerns the intellect, so that he cleanses ``from the defect
of nescience'' and perfects the intellect unto its end, which is known truth
(\latin{perfectio autem est consummatio in finem intellectus, qui est
veritas cognita}). Aquinas cites the \emph{Ecclesiastical
Hierarchy}: ``in the heavenly hierarchy the chastening of the inferior
essence is an enlightening of things unknown, that leads them to more
perfect knowledge'' (\emph{EH} 6, as quoted); and he gives an image from
bodily sight, ``cleansed by the removal of darkness; enlightened by the
diffusion of light; and perfected by being brought to the perception of the
colored object'' (ad~1). The Seraphim ``kindle'' others, then, by
persuasion to the love of God (ad~2), and Aristotle's higher appetite moving
the lower speaks of the sensitive appetite, which belongs to the same soul
and acts through an organ (ad~3). The commentary on the \emph{Sentences}
states the point about the Seraphim sharply: they are not called kindling
``as if they poured in the habit of charity, but because they enlighten
concerning the things that pertain to the love of charity'' (\latin{sed
quia illuminant de his quae ad amorem caritatis pertinent}; \emph{In II
Sent.} d.~9 q.~1 a.~4 ad~5).

Dionysius himself had joined the three acts to knowledge: ``the reception of
the supremely Divine Science is, both purification, and enlightenment, and
perfecting,---purifying, as it were, from ignorance, by the knowledge of
the more perfect revelations imparted to it according to fitness, and
enlightening by the self-same Divine knowledge \dots\ and perfecting further,
by the self-same Light, through the abiding science of the mysteries made
clearly manifest'' (\emph{CH} 7.3). Aquinas unfolds the sentence by the two
terms of every change. In receiving knowledge the starting point may be an
error contrary to the knowledge, or contrary dispositions, such as
impurity of soul or immoderate occupation with sensible things, or only the
absence of the knowledge, \latin{sicut cum in cognitione de die in diem
proficimus}; \latin{et sic tantummodo est accipere terminum a quo in
Angelis}. The goal is twofold: that by which the intellect is perfected to
know, whether form, light, or other medium, and the knowledge itself.
Purgation is the removal of nescience, illumination the first goal,
perfection the last; illumination and perfection differ \latin{sicut
formatio visus per speciem visibilis, et cognitio ipsius visibilis}, as the
forming of sight by the species of the visible differs from the knowing of
the visible itself. The ecclesiastical orders show the same
steps in their proper matter: the deacons had charge of the catechumens and
the possessed, in whom there were dispositions contrary to illumination;
the priests communicate and show the sacraments; the bishops opened to the
people the spiritual things veiled in the signification of the sacraments
(\emph{De veritate} q.~9 a.~3). Dionysius sums up the three ranks in the
same words: ``the Rank of the Hierarchs is consecrating and perfecting, that
of the Priests, illuminating and conducting to the light; and that of the
Leitourgoi purifying and discriminating'' (\emph{EH} 5.7).

\paragraph{The lower never enlightens the higher (a.~3).} The \emph{sed
contra} is Dionysius's ``Divine unalterable law, that inferior things are
led to God by the superior'' (\latin{hanc legem esse divinitatis
immobiliter firmatam, ut inferiora reducantur in Deum per superiora};
\emph{CH} 4; \emph{EH} 5). In Parker's text the Word of God teaches ``that
it came to us through Angels, as though the Divine regulation were laying
down this rule, that, through the first, the second are brought to the
Divine Being. For not only with regard to the superior and inferior minds,
but even for those of the same rank, this Law has been established by the
superessential supreme ordinance, that, within each Hierarchy, there are
first, and middle, and last ranks and powers, and that the more divine are
instructors and conductors of the less'' (\emph{CH} 4.3). Aquinas
determines: ``The inferior angels never enlighten the superior, but are
always enlightened by them.'' Order stands under order as cause under cause,
and it is fitting that something be done outside the order of a lower cause
for the sake of ordering it to a higher, ``as in human affairs the command
of the president is passed over from obedience to the prince.'' God works
miracles outside the order of corporeal nature so that men may be ordered to
knowledge of him. ``But the passing over of the order that belongs to
spiritual substances in no way belongs to the ordering of men to God; since
the angelic operations are not made known to us; as are the operations of
sensible bodies.'' The order of spiritual substances is therefore never
passed over by God, ``so that the inferiors are always moved by the
superior, and not conversely'' (\ST{I q.~106 a.~3 co.}).

The first objection draws on the Church: the ecclesiastical hierarchy
represents the heavenly, and the heavenly Jerusalem is ``our mother'' (Gal
4:26); yet in the Church the higher are taught by the lower, for ``you may
all prophesy, one by one, that all may learn and all may be exhorted'' (1
Cor 14:31). Aquinas answers that the ecclesiastical hierarchy imitates the
heavenly ``in some degree,'' not perfectly (\latin{aliqualiter, sed non
perfecte consequitur eius similitudinem}). In the heavenly hierarchy ``the
perfection of the order is in proportion to its nearness to God,'' and those
nearer to God are higher in grade and clearer in knowledge; in the Church
those who are nearer to God in holiness are sometimes in the lowest grade and
not eminent in learning, and some excel in one science and fail in another
(ad~1). The will of an angel turns freely to enlighten another, ``but the
angel's will is ever regulated by the Divine law which made the order in the
angels'' (ad~3).

Aquinas records in the disputed questions that the doctors had divided on
this point. Some held the order so firmly established that nothing ever
happens outside it; others held it established so as to be kept for the
most part, yet sometimes passed over for necessary causes, as the course of
nature is changed in miracles. He judges the first opinion more reasonable
for three reasons: it belongs to the dignity of the higher angels that the
lower be enlightened through them; \latin{quanto aliqua sunt Deo, qui est
summe immobilis, propinquiora, tanto debent esse magis immobilia}, so that the lower
bodies sometimes fail of their course while the heavenly bodies keep
theirs; and nothing in the state of nature is changed by divine power except
for something better, belonging to grace or glory, whereas \latin{statu
gloriae, in quo ordines Angelorum distinguuntur, nullus est altior status}
(\emph{De veritate} q.~9 a.~2).

\paragraph{The higher enlightens the lower in all that he knows (a.~4).}
The \emph{sed contra} gathers two authorities. The first Aquinas gives under
Gregory's name: ``In that heavenly country, though there are some excellent
gifts, yet nothing is held individually.'' The sentence is the Master's,
who writes \latin{In illa enim caelesti patria, ubi plenitudo boni est,
licet quaedam data sint excellenter, nihil tamen possidetur singulariter:
omnia enim in omnibus sunt, non quidem aequaliter, quia alii aliis sublimius
possident, quae tamen omnes habent} (Peter Lombard, \emph{Sent.} II d.~9
c.~3). Its source is Gregory's homily on the ten drachmas: \latin{Sic
quippe in illa summa civitate specialia quaedam singulorum sunt, ut tamen
sint communia omnium: et quod in se ex parte quisque habet, hoc in alio
ordine totum possidet} --- in that highest city certain things belong
specially to each, yet so as to be common to all; and what each has in part
in himself, he possesses whole in another order. Even what some have in a
way that others cannot, such as the names of Dominations and
Principalities, \latin{cuncta ibi singulorum sunt: quia per caritatem
Spiritus ab alio in aliis habentur} --- all things there belong to each,
because through the charity of the Spirit what one has is had by the others
in him (Gregory, \emph{Hom.\ in Evang.} 34.14). The second authority is
Dionysius, ``Each heavenly essence communicates to the inferior the gift
derived from the superior'' (\emph{CH} 15).

Aquinas answers from the nature of the good. ``Every creature participates
in the Divine goodness, so as to diffuse the good it possesses to others;
for it is of the nature of good to communicate itself to others''
(\latin{nam de ratione boni est quod se aliis communicet}). The more an
agent shares in the divine goodness, the more it strives to pass its
perfections to others; so Peter admonishes those who share it by grace, ``As
every man hath received grace, ministering the same one to another: as good
stewards of the manifold grace of God'' (1 Pet 4:10). ``Much more therefore
do the holy angels, who enjoy the plenitude of participation of the Divine
goodness, impart the same to those below them.'' Yet what is given is not
received as excellently below as above, ``and therefore the superior ever
remain in a higher order, and have a more perfect knowledge; as the master
understands the same thing better than the pupil who learns from him''
(\ST{I q.~106 a.~4 co.}). The higher's knowledge is called more universal
by reason of its more eminent mode (ad~1).

The second objection turns on the mystery of the Incarnation. The Master
holds that the higher angels knew it from the ages and the lower only when
it was accomplished, and Dionysius expounds the question of Psalm 23[24], in
which some angels ask as if not knowing, ``Who is this King of Glory?'' and
others answer as knowing, ``the Lord of hosts, he is the King of Glory'' (Ps
23[24]:10). In Parker's text the theologians show ``some of them as being
religiously instructed, by those of a higher rank, that He, Who was raised
to Heaven as Man, is Lord of the Heavenly Powers and King of Glory; and
others, as questioning Jesus Himself,'' and Dionysius marvels that even the
first ``do not ask directly, `Wherefore are Thy garments red?' but they
first raise the question among themselves, shewing that they desire to
learn'' (\emph{CH} 7.3; Isa 63:1--2). The Master sets the doctors side by
side. Jerome, on the Apostle's words that the manifold wisdom of God is made
known ``to the principalities and powers in heavenly places through the
church'' (Eph 3:10), held that the angelic dignities did not purely
understand the mystery until the Passion was complete and the preaching of
the Apostles spread; Augustine held that it was not hidden from them,
\latin{Illis ergo a saeculis innotuit supra memoratum mysterium}. Peter
Lombard notes the tension, \latin{Attende, lector, quia videntur dissentire
in hac sententia illustres doctores}, and, following Haymo, resolves it:
the angels of greater dignity, \latin{per quorum ministerium illa nuntiata
sunt}, knew it in part from the ages as familiars and messengers, the
lesser not until it was fulfilled and preached through the Church,
\latin{et tunc ab omnibus Angelis perfecte fuerunt cognita} (\emph{Sent.}
II d.~11 c.~2). Aquinas reads the Master gently: the lower angels were not wholly
ignorant of the mystery, ``but that they did not know it as fully as the
superior angels; and that they progressed in the knowledge of it afterwards
when the Mystery was accomplished'' (ad~2). His commentary on the same
distinction draws the finer line. The substance of the mystery, the
Incarnation and the Passion, all knew from the beginning; its conditions
and circumstances, \latin{quod sub tali praeside, vel tali hora}, they did
not, as the prophet announced the substance of the deed and the evangelist
recounts the manner of its fulfilment. By illumination the angels receive
nothing from men; but future contingents they know as they are fulfilled,
and so certain things about the mystery which they did not know they came to
know \latin{quando explebantur praedicantibus apostolis \dots\ non tamen ab
apostolis edocti}, when they were fulfilled as the Apostles preached, though
not taught by the Apostles (\emph{In II Sent.} d.~11 q.~2 a.~4). ``On whom the angels desire to look,'' says Peter of the
things now preached (1 Pet 1:12). The angels' knowledge of the mysteries
of grace is treated at \S\ref{sec:q57}.

The third objection fears that illumination would end once the lower knew
all the higher know. Aquinas answers with the history of the world: ``Till
the Judgment Day some new things are always being revealed by God to the
highest angels, concerning the course of the world, and especially the
salvation of the elect. Hence there is always something for the superior
angels to make known to the inferior'' (ad~3). ``There shall be joy before
the angels of God upon one sinner doing penance'' (Luke 15:10).

\angeldossier{\ST{I q.~106 aa.~1--4}; parallels \emph{De veritate} q.~9
aa.~1--3; \emph{In II Sent.} d.~9 q.~1 aa.~1--2, d.~11 q.~2 aa.~2,~4;
\emph{Summa contra gentiles} III.79--80; \ST{I q.~103 a.~6}.}
 {Common teaching that the higher angels enlighten the lower and send them
 to their ministries (a.~1; Dionysius, \emph{CH} 4.3, 7.3, 8.1; Gregory,
 \emph{Hom.\ in Evang.} 34.13; Isidore, \emph{Etym.} VII.5.8; Zach 2:3--4 [2:7--8]).
 Thomist position on the mode, by strengthening the lower intellect and by
 dividing the higher's universal conception (a.~1; \emph{De ver.} q.~9
 a.~1). Common teaching that God alone moves the will from within and that
 an angel moves it only by persuasion (a.~2); Thomist position that
 cleansing, enlightening, and perfecting among the angels all belong to the
 reception of knowledge (a.~2 ad~1; \emph{De ver.} q.~9 a.~3). Thomist
 position that the order of illumination is never passed over (a.~3; the
 contrary opinion is recorded in \emph{De ver.} q.~9 a.~2). Common teaching
 that the goods of heaven are held in common through charity, the higher
 possessing them more excellently (a.~4; Gregory, \emph{Hom.} 34.14;
 Lombard, \emph{Sent.} II d.~9 c.~3); Disputed among the Fathers how far
 the angels knew the mystery of the Incarnation before its fulfilment
 (Jerome and Augustine, as the Master reports them), with Aquinas's
 distinction of substance and circumstances (a.~4 ad~2; \emph{In II Sent.}
 d.~11 q.~2 a.~4).}
 {Dionysius, \emph{CH} 8 (\emph{sed contra} of a.~1), \emph{CH} 7 (a.~1 co.,
 a.~4 obj.~2), \emph{CH} 15 (a.~1 co., \emph{sed contra} of a.~4), \emph{CH}
 4 and \emph{EH} 5 (\emph{sed contra} of a.~3), \emph{CH} 12 (a.~4 obj.~1),
 \emph{EH} 6 (a.~2 ad~1), \emph{DN} 4 (a.~1 ad~1), all but \emph{EH} 6 read
 at the locus in Parker; Augustine, \emph{De div.\ qq.\ LXXXIII} q.~51 (a.~1
 obj.~3, read at the locus); Gregory, \emph{Hom.\ in Evang.} 34.13--14;
 Peter Lombard, \emph{Sent.} II d.~9 c.~3 and d.~11 c.~2 (\emph{sed contra}
 and obj.~2 of a.~4); Isidore, \emph{Etym.} VII.5.8; Eph 5:13, Eph 3:8--10,
 Jer 31:34, Exod 33:18--19, Gal 4:26, 1 Cor 14:31, 1 Pet 4:10, 1 Pet 1:12,
 Ps 23[24]:8--10, Isa 63:1--2, Luke 15:10; Aristotle, \emph{De anima} III (a.~2
 obj.~3, cited by Aquinas).}
 {The opinion, recorded and judged less reasonable in \emph{De ver.} q.~9
 a.~2, that the order of illumination is kept for the most part but
 sometimes passed over, as the course of nature is in miracles (a.~3);
 Jerome, as reported by Peter Lombard, that the angels did not understand
 the mystery of the Incarnation until the Passion and the preaching of the
 Apostles (a.~4 obj.~2).}

\subsection{The Speech of the Angels (\ST{I q.~107})}\label{sec:q107}

The five articles ask whether one angel speaks to another; whether the
lower speaks to the higher; whether an angel speaks to God; whether local
distance has any effect on angelic speech; and whether all the angels know
what one says to another (\ST{I q.~107 pr.}). The Apostle names ``the
tongues of men and of angels'' (1 Cor 13:1), and Gregory gives the rule for
understanding them. The Lord who is ``the supreme and unbounded Spirit'' and
Satan, ``who is invested with no fleshly nature,'' do not draw breath into
the lungs and give it back through the throat in articulate voice; ``but
when the Incomprehensible Nature speaks to an invisible nature, it behoves
that our imagination rising above the properties of our corporeal speech
should be lifted to the sublime and unknown methods of interior speech''
(Gregory, \emph{Moralia} II.8). Aquinas opens the question with the same
words: \latin{dignum est ut mens nostra, qualitatem corporeae locutionis
excedens, ad sublimes atque incognitos modos locutionis intimae suspendatur}
(\ST{I q.~107 a.~1 co.}). And Gregory adds at once that such speech ``is by
no means characterized by one and the same form. For it is after one method
that God speaks to the Angels, and after another that the Angels speak to
God; in one manner that God speaks to the souls of Saints, in another that
the souls of Saints speak to God; in one way God speaks to the devil, in
another the devil speaks to God'' (\emph{Moralia} II.8). Aquinas's five
articles are a disciplined unfolding of the first two of these modes and of
the speech of angel with angel.

\paragraph{One angel speaks to another (a.~1).} The \emph{sed contra} is
the Apostle, ``If I speak with the tongues of men and of angels'' (1 Cor
13:1). The account rests on the will's command over the intellect. An
intelligible object is present to the intellect in three ways: habitually,
``or in the memory, as Augustine says''; as actually considered or
conceived; and as related to something else. It passes from the first state
to the second by the command of the will, ``and hence in the definition of
habit these words occur, `which anyone uses when he wills.''' It passes from
the second to the third likewise by the will, ``for by the will the concept
of the mind is ordered to something else, as, for instance, either to the
performing of an action, or to being made known to another.'' When the mind
turns to consider what it holds in habit, one speaks to oneself, ``for the
concept of the mind is called `the interior word.''' And when the concept of
an angelic mind is ordered by the angel's own will to be made known to
another, it becomes known to the other; ``and in this way one angel speaks
to another; for to speak to another only means to make known the mental
concept to another'' (\latin{Nihil est enim aliud loqui ad alterum, quam
conceptum mentis alteri manifestare}) (\ST{I q.~107 a.~1 co.}).

The first objection is Gregory's own: in the resurrection the body will not
hide one mind from another, much less is one angel's mind hidden from
another. Gregory's page is a vision of the city of glass. Job's ``The gold
and the glass cannot equal it'' is read of ``that heavenly Country, that
society of blessed citizens, whose hearts mutually one with another at once
shine with brightness, and are transparent by pureness,'' the city John saw
``of pure gold like unto clear glass'' (Apoc 21:18). ``For there the mind of
every person no bodily frame of limbs will hide from the eyes of his fellow,
but the interior will be given to view'' (\emph{Moralia} XVIII.77--78).
Aquinas answers that our concept is shut in by a twofold obstacle. The first
is the will, which can keep the concept within or direct it outward, and in
this respect ``God alone can see the mind of another, according to 1
Corinthians 2:11: `What man knoweth the things of a man, but the spirit of a
man that is in him?''' The second is the body, so that even when the will
directs the concept outward some sensible sign must be used. Gregory again
supplies the image: ``we, that we may express outwardly the things which we
are inwardly sensible of, deliver these through the organ of the throat, by
the sounds of the voice, since to the eyes of others we stand as it were
behind the partition of the body, within the secret dwelling place of the
mind; but when we desire to make ourselves manifest, we go forth as though
through the door of the tongue, that we may shew what kind of persons we are
within'' (\emph{Moralia} II.8). ``But an angel is under no such obstacle,''
and \latin{quam cito vult manifestare suum conceptum, statim alius
cognoscit} (ad~1). The glass of the heavenly city, then, is the removal of
the second obstacle, not the first: the will's own gate remains, and the
angels open it to one another.

Exterior speech by voice is necessary to us because of the body, and does
not befit an angel; his speech is interior, ``and this includes not only
the interior speech by mental concept, but also its being ordered to
another's knowledge by the will. So the tongue of an angel is called
metaphorically the angel's power, whereby he manifests his mental concept''
(ad~2). There is no need to arouse the attention of the good angels,
``inasmuch as they always see each other in the Word''; but since by
nature the angels could
speak to one another, ``and even now the bad angels speak to each other,''
Aquinas adds that as sense is aroused by a sensible sign, so the mind of an
angel can be aroused to attention ``by some intelligible power'' (ad~3).

The disputed questions give the same account through the three states of a
natural form: imperfect, as in forms still coming to be; perfect in itself;
and perfect so as to communicate its perfection to another, \latin{aliquid
enim est in se lucidum, quod alia illuminare non potest}. The intelligible
form is in the intellect as in habit, as in act for the one who thinks, and
as ordered to another, and the will makes each transition; \latin{Et quando
sic est, tunc alius Angelus eius cognitionem percipit; et secundum hoc
dicitur alteri Angelo loqui}. Such speech is required because an angel does
not know the secrets of another's heart directly (\S\ref{sec:q57}); and so it would be with us, \latin{si
intellectus noster posset ferri in intelligibilia immediate} (\emph{De
veritate} q.~9 a.~4). The commentary on
the \emph{Sentences} had described the matter otherwise. There the concept
ordered to another is called the word of the heart, and speech is its
coordination with something the other angel can naturally see in the
speaker, which becomes a sign expressive of the interior concept: \latin{et
talis expressio vocatur locutio, non quidem vocalis, sed intellectualibus
signis expressa; et virtus exprimendi dicitur lingua eorum} (\emph{In II
Sent.} d.~11 q.~2 a.~3). The \emph{Summa} and the disputed questions need no
intermediate sign; the will's ordering of the concept suffices.

\paragraph{The lower speaks to the higher (a.~2).} The \emph{sed contra} is
again Dionysius's exposition of the Psalm, in which the lower angels asked
the higher, ``Who is this King of Glory?'' (\emph{CH} 7.3). The article
distinguishes speech from illumination. ``Every angelic enlightening is an
angelic speech; but on the other hand, not every speech is an enlightening.''
What the mind conceives can be referred to two principles: to God, the first
truth, and to the will of the one who understands. ``Because truth is the
light of the intellect, and God Himself is the rule of all truth,'' the
manifestation of what depends on the first truth is both speech and
enlightenment, as when one man says to another, ``Heaven was created by
God.'' The manifestation of what depends on the will is speech only, as when
one says, ``I wish to learn this; I wish to do this or that.'' The reason is
that ``the created will is not a light, nor a rule of truth; but
participates of light. Hence to communicate what comes from the created will
is not, as such, an enlightening. For to know what you may will, or what you
may understand does not belong to the perfection of my intellect; but only
to know the truth in reality.'' Higher and lower are said by relation to
God; but ``as regards the will as the principle, he who wills is first and
supreme''; and so in this second kind of speech the higher speak to the
lower and the lower to the higher (\ST{I q.~107 a.~2 co.}). Every speech of
God to the angels, however, is an enlightening, ``because since the will of
God is the rule of truth, it belongs to the perfection and enlightenment of
the created mind to know even what God wills'' (ad~3).

Gregory's description of that divine speech is among the loveliest pages of
the \emph{Moralia}. ``For because no corporeal obstacle is in the way of a
spiritual being, God speaks to the holy Angels in the very act of His
revealing to their hearts His inscrutable secrets, that whatsoever they
ought to do they may read it in the simple contemplation of truth, and that
the very delights of contemplation should be like a kind of vocal precepts,
for that is as it were spoken to them as hearers which is inspired into
them as beholders'' (\emph{Moralia} II.9). When God said ``Come, let us go
down, and there confound their language,'' he spoke to those close about
him: ``never to depart from Him in heart, is as it were to be always coming
to Him by a kind of steady motion \dots\ The Angels ascend in that they
behold their Creator; the Angels descend in that by a strict examination
they put down that which exalts itself in unlawful measure'' (ibid.).

The disputed questions give the distinction its sharpest form. A mind fails
to know something either because the thing is absent, as the deeds of past
times or distant places, or because the mind is not strong enough to reach
what it already holds, as all conclusions are held in first principles.
\latin{Locutio igitur proprie est qua aliquis ducitur in cognitionem ignoti,
per hoc quod fit sibi praesens quod alias erat ei absens}; illumination is
the strengthening of the intellect to know beyond what it knew. There can
be speech without illumination, as when histories are recited to me; but
illumination always has speech joined to it, since the higher angel, whose
knowledge is by more universal forms, must conceive in himself what he would
teach in a mode the lower can grasp, and manifest that concept to him. Hence
\latin{illa locutione quae illuminationi adiungitur, superiores solum
inferioribus, loquuntur; sed secundum aliam locutionem indifferenter
loquuntur et superiores inferioribus, et e converso} (\emph{De veritate}
q.~9 a.~5).

\paragraph{An angel speaks to God (a.~3).} The \emph{sed contra} is the
intercession of Zechariah's vision: ``The angel of the Lord answered, and
said: O Lord of hosts, how long wilt thou not have mercy on Jerusalem'' (Zach
1:12). One thing is ordered to another in two ways. It may be ordered in
order to give something, as agent to patient and teacher to learner, and in
this way an angel never speaks to God, ``either of what concerns the truth,
or of whatever depends on the created will; because God is the principle and
source of all truth and of all will.'' Or it may be ordered in order to
receive, as the passive to the agent and the disciple to the master, ``and in
this way an angel speaks to God, either by consulting the Divine will of what
ought to be done, or by admiring the Divine excellence which he can never
comprehend''; ``thus Gregory says (Moral. ii) that `the angels speak to God,
when by contemplating what is above themselves they rise to emotions of
admiration''' (\ST{I q.~107 a.~3 co.}). Speech is not always ordered to
making something known to another; sometimes it is ordered to having
something made known to the speaker, ``as when the disciples ask instruction
from the master'' (ad~1). ``The angels are ever speaking to God in the sense
of praising and admiring Him and His works; but they speak to Him by
consulting Him about what ought to be done whenever they have to perform any
new work, concerning which they desire enlightenment'' (ad~2).

Gregory's text is a meditation on the heavenly liturgy. ``It is after another
manner that the Angels speak to God, as in the Revelation of John also they
say, Worthy is the Lamb that was slain to receive power, and riches, and
wisdom; for the voice of the Angels in the praises of God is the very
admiration itself of inward contemplation. To be struck dumb at the marvels
of Divine goodness is to utter a voice, for the emotion of the heart excited
with a feeling of awe is a mighty utterance of voice to the ears of a Spirit
that is not circumscribed. This voice unfolds itself as it were in distinct
words, while it moulds itself in the innumerable modes of admiration. God
then speaks to the Angels when His inner will is revealed to them as the
object of their perception; but the Angels speak to the Lord when by means of
this, which they contemplate above themselves, they rise to emotions of
admiration'' (\emph{Moralia} II.10; Apoc 5:12). The same speech, turned
toward what is to be done, is consultation. When the angel of Daniel says
that ``the prince of the kingdom of the Persians resisted me one and twenty
days'' (Dan 10:13), Aquinas follows Gregory in holding both to be good
angels: the depth of God's judgments exceeds the angels, the merits of the
nations they have in charge are contrary, and each refers the merits of his
own to the divine knowledge, awaiting the sentence. The contrary relation is
called a combat, \latin{et eorum concordia est in divinae illuminationis
perceptione, per quam de divina voluntate instruuntur: hoc enim omnes
concorditer volunt quod percipiunt Deum velle} (\emph{In II Sent.} d.~11
q.~2 a.~5; the angels of the nations are treated at \S\ref{sec:q113}).
``Power and terror are with him, who maketh peace in his high places'' (Job
25:2).

\paragraph{Local distance does not affect angelic speech (a.~4).} The
\emph{sed contra} is the rich man in hell, who spoke to Abraham across the
great gulf (Luke 16:24); much less can distance hinder one angel's speech to
another. ``The angelic speech consists in an intellectual operation,'' and
the intellectual operation of an angel abstracts from ``here and now,'' as
even ours does, except accidentally by reason of phantasms, ``which do not
exist at all in an angel.'' In what is wholly abstracted from place and
time, neither difference of time nor distance of place has any effect
(\ST{I q.~107 a.~4 co.}). The first objection quotes Damascene, ``An angel
works where he is.'' Damascene's chapter on place says that the angel,
``although not contained in place with figured form as is body, yet is
spoken of as being in place because he has a mental presence and energises
in accordance with his nature, and is not elsewhere but has his mental
limitations there where he energises'' (\emph{De fide orthodoxa} I.13).
Aquinas grants that speech, being interior, is in the angel who speaks and so
where he is; but ``as local distance does not prevent one angel seeing
another, so neither does it prevent an angel perceiving what is ordered to
him on the part of another; and this is to perceive his speech'' (ad~1).
In the disputed questions he restricts Damascene's axiom to the operation
by which an angel acts on a body, which alone is \latin{situalis} (\emph{De
veritate} q.~9 a.~6 ad~2; on angelic place, \S\ref{sec:q52}).

The second objection is the Seraphim who ``cried one to another'' (Isa 6:3).
The cry, Aquinas answers, ``is not a bodily voice raised by reason of the
local distance; but is taken to signify the magnitude of what is said, or
the intensity of the affection, according to what Gregory says (Moral. ii):
`The less one desires, the less one cries out''' (ad~2). Gregory's sentence
stands in his exposition of the souls under the altar who ``cried with a loud
voice'' (Apoc 6:10): ``For their great cry is their great longing; for every
one cries the less, the less he desires; and he utters the louder voice in
the ears of an uncircumscribed Spirit in proportion as he more entirely pours
himself out in desire of Him, and so the words of souls are their very
desires'' (\emph{Moralia} II.11). Dionysius reads the same cry of the
Seraphim as the law of illumination within one order: ``the theologians say
that the most Divine Seraphim cry one to another, indicating distinctly, as I
think by this, that the first impart their knowledge of divine things to the
second'' (\emph{CH} 10.2). Isidore finds in the thrice-repeated ``Holy'' the
mystery of the Trinity: \latin{Quod clamat ter Sanctus alter ad alterum,
Trinitatis in una divinitate demonstrat mysterium} (\emph{Etym.} VII.5.33).

\paragraph{One angel may speak to another alone (a.~5).} The \emph{sed
contra} is from our own experience: ``One man can speak to another alone;
much more can this be the case among the angels.'' Since the concept of one
angel is perceived by another because its possessor orders it by his will
to the other, and a thing can be ordered to one and not to another, ``the
concept of one (angel) may be known by one and not by another''; and so one
angel can perceive the speech of another while others do not, ``not through
the obstacle of local distance, but on account of the will so ordering''
(\ST{I q.~107 a.~5 co.}). The third objection returns to Dionysius's law
that each heavenly being communicates to the others what it learns (\emph{CH}
15). Aquinas answers with the two principles of a.~2. ``Enlightenment is of
those truths that emanate from the first rule of truth, which is the
principle common to all the angels; and in that way all enlightenments are
common to all. But speech may be of something ordered to the principle of the
created will, which is proper to each angel; and in this way it is not
necessary that these speeches should be common to all'' (ad~3). The disputed
questions add that if an angel by his own will becomes actual in some
species in order to one angel only, his speech is perceived by that one
alone; \latin{si vero in ordine ad plures, percipietur a pluribus}
(\emph{De veritate} q.~9 a.~7). The communion of the angels in truth is
complete; the secrecy of the created will is respected; and the one Will
before whom every heart is open is God's.

\angeldossier{\ST{I q.~107 aa.~1--5}; parallels \emph{De veritate} q.~9
aa.~4--7; \emph{In II Sent.} d.~11 q.~2 aa.~3,~5.}
 {Common teaching that the angels speak to one another and to God (aa.~1,
 3; 1 Cor 13:1; Zach 1:12; Isa 6:3; Apoc 5:11--12; Gregory, \emph{Moralia}
 II.8--10). Thomist position on the account of angelic speech as the
 ordering of the mental concept to another by the will, without
 intermediate sign (a.~1; \emph{De ver.} q.~9 a.~4; the commentary on the
 \emph{Sentences} places the speech in intellectual signs), on the
 distinction of speech from illumination by the two principles of truth and
 will (aa.~2,~5), on the indifference of angelic speech to local distance
 (a.~4), and on the freedom of the speaker to be heard by one alone (a.~5).
 Common teaching, with Gregory, that the angels' speech to God is praise
 and admiration and consultation about what is to be done (a.~3).}
 {1 Cor 13:1 (\emph{sed contra} of a.~1); Dionysius, \emph{CH} 7 (\emph{sed
 contra} of a.~2), \emph{CH} 15 (a.~5 obj.~3), and \emph{CH} 10.2, read at
 the locus in Parker; Zach 1:12 (\emph{sed contra} of a.~3); Luke 16:24
 (\emph{sed contra} of a.~4); Isa 6:3 (a.~4 obj.~2); Gregory,
 \emph{Moralia} II.8--11 (a.~1 co.\ and ad~1, a.~2 obj.~3, a.~3 co., a.~4
 ad~2) and XVIII.77--78 (a.~1 obj.~1), read at the locus in the Library of
 the Fathers translation; Damascene, \emph{De fide orthodoxa} I.13 (a.~4
 obj.~1, read at the locus); Augustine, \emph{De Trinitate} XIV (a.~1 co.,
 cited by Aquinas); 1 Cor 2:11 (a.~1 ad~1); Apoc 21:18; Dan 10:13; Isidore,
 \emph{Etym.} VII.5.33.}
 {The gloss on 1 Cor 13:1 that the speeches of the angels are the
 illuminations by which the higher enlighten the lower, so that the lower
 do not speak to the higher (a.~2 obj.~1); Damascene's axiom that an angel
 works where he is, taken to limit speech by distance (a.~4 obj.~1,
 restricted by Aquinas to action on bodies).}
